\documentclass[a4paper,14pt]{extarticle}

% Настройки кодировки, шрифтов и локализации
\usepackage[utf8]{inputenc}
\usepackage[T2A]{fontenc}
\usepackage[russian]{babel}
\usepackage{times} 

% Геометрия страницы по ГОСТ (Левое поле — 30 мм, правое — 15 мм, верхнее/нижнее — 20 мм)
\usepackage[left=30mm, right=15mm, top=20mm, bottom=20mm]{geometry}

% Набор пакетов для оформления графики, кода и математики
\usepackage{graphicx}
\usepackage{hyperref}
\usepackage{listings}
\usepackage{color}
\usepackage{caption}
\usepackage{amsmath}
\usepackage{indentfirst} 

% Конфигурация интерактивных ссылок в документе
\hypersetup{
	colorlinks=true,
	linkcolor=black,
	filecolor=black,      
	urlcolor=blue,
	citecolor=black
}

% Цветовая палитра для листингов кода (исключительно безопасные тона)
\definecolor{codegray}{rgb}{0.5,0.5,0.5}
\definecolor{backcolour}{rgb}{0.98,0.98,0.96}
\definecolor{keywordblue}{rgb}{0.05,0.2,0.65}
\definecolor{stringgreen}{rgb}{0.1,0.5,0.1}

\lstset{
	backgroundcolor=\color{backcolour},   
	commentstyle=\color{codegray}\itshape,
	keywordstyle=\color{keywordblue}\bfseries,
	stringstyle=\color{stringgreen},
	basicstyle=\ttfamily\footnotesize,
	breakatwhitespace=false,         
	breaklines=true,                 
	captionpos=b,                    
	keepspaces=true,                 
	numbers=left,                    
	numbersep=7pt,                  
	showspaces=false,                
	showstringspaces=false,
	showtabs=false,                  
	tabsize=2,
	frame=single,
	rulecolor=\color{codegray}
}

\begin{document}
	
	% ================= ТИТУЛЬНЫЙ ЛИСТ =================
	\begin{titlepage}
		\begin{center}
			\textsc{\large Министерство науки и высшего образования Российской Федерации}\\[0.4cm]
			\textsc{\large Южный федеральный университет}\\[0.4cm] %[cite: 1]
			\textsc{\large Физический факультет}\\[4.5cm] %[cite: 1]
			
			\textbf{\Large КУРСОВОЙ ПРОЕКТ}\\[0.5cm]
			\large по дисциплине «Администрирование ОС Linux»\\[1cm]
			
			\textbf{\Large Тема: Развертывание и настройка NoSQL базы данных MongoDB в режиме Replica Set}\\[3.5cm]
		\end{center}
		
		\begin{flushright}
			\large
			\textbf{Исполнители (Группа):} \\
			\textbf{Мизиев Р. С.} \\ %[cite: 1]
			\textbf{Сорокина Д.Б.} \\
			\textbf{Деккер П.И.} \\[0.8cm]
			\textbf{Проверил:} \\
			Преподаватель кафедры \\
		\end{flushright}
		
		\vfill
		\begin{center}
			\large Ростов-на-Дону, \the\year
		\end{center}
	\end{titlepage}
	
	% ================= СОДЕРЖАНИЕ =================
	\newpage
	\tableofcontents
	\newpage
	
	% ================= ВВЕДЕНИЕ =================
	\section*{Введение}
	\addcontentsline{toc}{section}{Введение}
	Современный этап развития информационных технологий характеризуется лавинообразным ростом объемов генерируемых данных и ужесточением требований к доступности сетевых сервисов. В условиях парадигмы высокой доступности (High Availability) традиционная одиночная инсталляция сервера баз данных превращается в критический вектор отказа всей системы (Single Point of Failure). Любой аппаратный сбой, деградация сетевого канала или некорректное распределение оперативной памяти на хост-машине могут повлечь за собой полную остановку бизнес-логики приложений.
	
	Для преодоления ограничений классических реляционных СУБД в контексте горизонтального масштабирования все более широкое распространение находят решения класса NoSQL. Документоориентированная СУБД MongoDB занимает лидирующие позиции в данном сегменте благодаря гибкой схеме документов (BSON), высокой производительности и нативной поддержке механизмов распределения данных. 
	
	Однако надежность и производительность распределенной системы напрямую зависят от качества администрирования базовой операционной системы. Развертывание таких комплексов вручную неизбежно приводит к человеческим ошибкам и «дрейфу конфигураций» (configuration drift). В связи с этим актуальным становится применение методологии «Инфраструктура как код» (Infrastructure as Code, IaC), позволяющей декларативно описывать, разворачивать и контролировать состояние вычислительной среды с помощью автоматизированных сценариев ОС Linux.
	
	\textbf{Цель данного проекта} — комплексное проектирование, развертывание, системная оптимизация и многоуровневая защита распределенного кластера NoSQL СУБД MongoDB в режиме репликации с использованием современных средств автоматизации Linux-окружения.
	
	\newpage
	
	% ================= РАЗДЕЛ 1: РОЛИ =================
	\section{Распределение ролей в группе и индивидуальный вклад}
	Разработка, тестирование и обеспечение безопасности распределенного кластера СУБД выполнялись проектной группой из трех человек с четким разграничением обязанностей и зон инженерной ответственности. 
	
	\subsection{Системный администратор — Мизиев Р. С.}
	В зону ответственности системного администратора входили задачи по архитектурному планированию и конфигурации системных компонентов распределенной базы данных. Были выполнены следующие этапы:
	\begin{itemize}
		\item Проектирование топологии кластера и определение весов и приоритетов выборности узлов;
		\item Разработка главного архитектурного манифеста распределенного окружения;
		\item Профилирование и тонкая настройка параметров движка хранения данных, расчет допустимых лимитов оперативной памяти во избежание сбоев подсистемы виртуальной памяти ядра Linux;
		\item Оптимизация конфигурационных файлов СУБД для обеспечения бесконфликтной работы нескольких изолированных демонов на едином вычислительном узле.
	\end{itemize}
	
	\subsection{DevOps-инженер — Сорокина Д.Б.}
	DevOps-инженер сфокусировался на автоматизации жизненного цикла стенда, обеспечении воспроизводимости окружения и непрерывной валидации кода:
	\begin{itemize}
		\item Разработка bash-сценариев автоматической инициализации окружения, генерации структуры каталогов и подготовки переменных сред;
		\item Написание скриптов автоматизации процессов «горячего» резервного копирования данных с последующей очисткой устаревших архивов по алгоритмам ротации;
		\item Построение сквозных интеграционных тестов для проверки сетевой и транзакционной готовности СУБД;
		\item Создание конвейера непрерывной интеграции на удаленных ранерах для автоматического синтаксического анализа и контроля качества кода (Linting).
	\end{itemize}
	
	\subsection{Специалист по информационной безопасности — Деккер П.И.}
	Специалист по информационной безопасности обеспечивал защиту конфиденциальности, целостности и доступности данных на всех уровнях абстракции системы:
	\begin{itemize}
		\item Внедрение механизмов криптографической межсерверной аутентификации нод кластера с ограничением прав доступа на уровне файловой системы Linux;
		\item Разработка и автоматизация скрипта защиты хоста посредством конфигурации системного брандмауэра UFW по концепции минимизации площади атаки;
		\item Проектирование изолированных виртуальных сетей для исключения возможности несанкционированного межконтейнерного перехвата трафика базы данных;
		\item Включение механизмов строгой ролевой модели авторизации доступа (RBAC) внутри СУБД.
	\end{itemize}
	
	\newpage
	
	% ================= РАЗДЕЛ 2: ТЕОРИЯ =================
	\section{Теоретические основы распределенных NoSQL систем}
	При проектировании распределенных хранилищ данных ключевым теоретическим базисом выступает теорема CAP, сформулированная Эриком Брюером. Согласно данной теореме, распределенная система способна одновременно обеспечить не более двух из трех следующих свойств:
	\begin{itemize}
		\item \textbf{C (Consistency)} — согласованность данных. Все узлы в один и тот же момент времени видят идентичные данные;
		\item \textbf{A (Availability)} — доступность. Каждый запрос к работающему узлу завершается корректным ответом;
		\item \textbf{P (Partition Tolerance)} — устойчивость к разделению. Система продолжает функционировать при потере связи между узлами.
	\end{itemize}
	
	В условиях реальных физических сетей возникновение задержек и разрывов каналов связи неизбежно, поэтому выбор архитектуры всегда балансирует между парадигмами CP и AP. Кластер MongoDB в режиме репликации (Replica Set) является классической **CP-системой**. В случае сетевого расщепления кластер жертвует доступностью на запись ради сохранения строгой согласованности данных и предотвращения рассинхронизации.
	
	Расширенная теорема PACELC дополняет этот баланс: если система работает нормально (аббревиатура Else), разработчик должен выбрать между задержкой (Latency, L) и согласованностью (Consistency, C). Репликация в MongoDB позволяет гибко управлять этим компромиссом на уровне отдельных запросов за счет параметров \texttt{Read Concern} (определяет порог актуальности читаемых данных) и \texttt{Write Concern} (задает количество узлов, обязанных подтвердить успешную запись до возврата ответа клиенту).
	
	Механизм репликации опирается на концепцию журнала операций \textbf{Oplog} (Operations Log). Данный журнал представляет собой ограниченную в размерах коллекцию, в которую Primary-узел записывает все модификации данных. Вторичные узлы непрерывно реплицируют Oplog лидера и применяют содержащиеся в нем инструкции к своим локальным копиям данных. Операции в Oplog являются \textit{идемпотентными}: их повторное выполнение на Secondary-нодах гарантированно приводит к одному и тому же результату, предотвращая повреждение структуры БД при сетевых сбоях переповтора пакетов.
	
	\newpage
	
	% ================= РАЗДЕЛ 3: АРХИТЕКТУРА =================
	\section{Проектирование архитектуры и сетевая изоляция кластера}
	Для минимизации аппаратных затрат при сохранении полноценного функционала отказоустойчивости в рамках проекта была выбрана классическая трехнодовая топология \textbf{Primary-Secondary-Arbiter (PSA)}. 
	
	Схема функционирования ролей узлов выглядит следующим образом:
	\begin{enumerate}
		\item \textbf{mongo-primary}: Единственная точка входа для транзакций записи. Узел координирует состояние кластера, принимает изменения от клиентов и генерирует Oplog. Приоритет узла в конфигурации равен 2.
		\item \textbf{mongo-secondary}: Узел-реплика. Постоянно синхронизирует состояние с лидером. Обладает приоритетом 1, что позволяет ему автоматически инициировать процедуру выборов и занять роль Primary в случае недоступности текущего лидера.
		\item \textbf{mongo-arbiter}: Специализированный технический узел. Он принципиально не имеет дискового хранилища под данные СУБД, не участвует в процессах репликации и не может стать лидером. Его единственное назначение — обеспечение нечетного числа голосов при голосовании. Это предотвращает ситуацию «равенства голосов» и защищает кластер от логического тупика.
	\end{enumerate}
	
	Сетевая безопасность спроектирована по принципу полной изоляции данных. С помощью стандартных средств контейнеризации Linux развернута обособленная виртуальная сеть \texttt{backend-net} с типом драйвера \texttt{bridge}. Внутри данной сети каждому контейнеру выделяется зафиксированный локальный IP-адрес и DNS-имя. 
	
	Внешний сетевой интерфейс хост-машины полностью закрыт для трафика СУБД на порту \texttt{27017}. Для администратора развернута изолированная веб-панель \texttt{Mongo Express}, проксирующая запросы к базе данных через внутренние каналы виртуальной сети и предоставляющая наружу защищенный веб-интерфейс.
	
	\newpage
	
	% ================= РАЗДЕЛ 4: IAC =================
	\section{Автоматизация развертывания и концепция Infrastructure as Code}
	Развертывание распределенной инфраструктуры реализовано в рамках парадигмы IaC с помощью декларативного инструмента Docker Compose. Это позволяет полностью исключить дрейф конфигураций и гарантировать идентичность сред разработки, тестирования и промышленной эксплуатации.
	
	\subsection{Декларативное описание стека}
	Ниже приведен листинг конфигурационного файла, описывающий параметры инициализации первичного узла кластера:
	\begin{lstlisting}
		services:
		mongo-primary:
		image: mongo:6.0.28
		container_name: mongo-primary
		restart: always
		command: [
		"mongod", 
		"--bind_ip_all", 
		"--replSet", "${MONGO_REPLICA_SET_NAME}", 
		"--keyFile", "/opt/keyfile/replica.key"
		]
		environment:
		MONGO_INITDB_ROOT_USERNAME: ${MONGO_INITDB_ROOT_USERNAME}
		MONGO_INITDB_ROOT_PASSWORD: ${MONGO_INITDB_ROOT_PASSWORD}
		ports:
		- "27017:27017"
		volumes:
		- mongo_primary_data:/data/db
		- ../config/mongodb/mongod.conf:/etc/mongod.conf:ro
		- ../deploy/docker/mongodb/keys/replica.key:/opt/keyfile/replica.key:ro
		networks:
		- backend-net
	\end{lstlisting}
	
	\subsection{Автоматизация жизненного цикла окружения}
	Сценарий \texttt{bootstrap.sh} берет на себя выполнение рутинных низкоуровневых операций в ОС Linux. Он выполняет валидацию наличия конфигурационного файла секретов \texttt{.env}, инициализирует файловые каталоги на хосте и генерирует общий секретный ключ авторизации с помощью криптографических библиотек ядра:
	\begin{lstlisting}[language=bash]
		if [[ ! -f "${KEY_FILE}" ]]; then
		openssl rand -base64 756 > "${KEY_FILE}"
		chmod 400 "${KEY_FILE}"
		fi
	\end{lstlisting}
	Управление развернутым стендом делегировано утилите \texttt{make}. В файле \texttt{Makefile} прописаны абстрактные цели для сборки окружения, запуска интеграционных тестов и очистки дискового пространства хоста.
	
	\newpage
	
	% ================= РАЗДЕЛ 5: OPTIMIZATION =================
	\section{Системная оптимизация и комплексная безопасность Linux-окружения}
	
	\subsection{Профилирование оперативной памяти и оптимизация WiredTiger}
	Одной из наиболее распространенных причин падения баз данных в контейнерах под управлением Linux является некорректное управление подсистемой кэширования. Движок хранения данных WiredTiger по умолчанию использует алгоритм, резервирующий под свои нужды до 50\% доступной оперативной памяти хост-компьютера за вычетом 1 Гигабайта. 
	
	При параллельном запуске трех независимых экземпляров СУБД на одном сервере их суммарные требования к оперативной памяти мгновенно приводят к исчерпанию физического объема ОЗУ и swap-пространства. В этот момент ядро операционной системы Linux активирует подсистему Out-Of-Memory (OOM) Killer. По алгоритму подсчета штрафных очков `badness score` именно процессы СУБД уничтожаются ядром в первую очередь для предотвращения падения всей ОС.
	
	Для стабилизации системы в конфигурационном файле \texttt{mongod.conf} размер внутреннего кэша был оптимизирован и жестко зафиксирован:
	\begin{lstlisting}
		storage:
		dbPath: /data/db
		engine: wiredTiger
		wiredTiger:
		engineConfig:
		cacheSizeGB: 1.0
	\end{lstlisting}
	Данная конфигурация гарантирует предсказуемый профиль потребления ресурсов оперативной памяти всего стенда.
	
	\subsection{Укрепление безопасности операционной системы (Hardening)}
	Многоуровневый периметр защиты включает в себя жесткое разграничение прав доступа в файловой системе POSIX. Файл криптографического ключа \texttt{replica.key} защищается маской прав \texttt{chmod 400}. Это лишает права на чтение и запись любые учетные записи операционной системы, кроме непосредственного владельца процесса.
	
	На уровне сетевого стека ядра Linux автоматизировано применение брандмауэра \texttt{UFW} (Uncomplicated Firewall). Сценарий \texttt{hardening.sh} переводит межсетевой экран в режим блокирования входящего трафика по умолчанию. Для внешней сети открываются исключительно строго регламентированные порты: \texttt{22} для защищенного удаленного администрирования по протоколу SSH, а также веб-интерфейсы. Прямой доступ к портам \texttt{27017} узлов СУБД из внешних сетей полностью заблокирован.
	
	\newpage
	
	% ================= РАЗДЕЛ 6: MAINTENANCE =================
	\section{Обслуживание, интеграционное тестирование и CI/CD}
	
	\subsection{Стратегия «горячего» резервного копирования}
	Для обеспечения непрерывности процессов управления данными разработан скрипт \texttt{backup.sh}, выполняющий резервное копирование без остановки СУБД. Процесс опирается на штатную утилиту \texttt{mongodump}, которая запускается внутри контейнера `mongo-primary`, считывает снапшот данных и выполняет его поточное сжатие с применением алгоритма GZIP. 
	
	Для предотвращения деградации дисковой подсистемы хоста из-за накопления старых архивов, в сценарий внедрен механизм автоматической ротации:
	\begin{lstlisting}[language=bash]
		find "${BACKUP_DIR}" -type f -name "*.gz" -mtime +7 -exec rm {} \;
	\end{lstlisting}
	
	\subsection{Автоматизация тестирования и конвейер CI/CD}
	Валидация работоспособности развернутой репликации осуществляется интеграционным скриптом \texttt{test\_mongo.sh}. Он выполняет циклическую проверку доступности сетевого сокета, опрашивает статус инициализации кластера через команду \texttt{rs.status().ok} и осуществляет тестовую запись документа в коллекцию.
	
	Весь проект интегрирован с облачным репозиторием и управляется конвейером автоматических проверок GitHub Actions на базе манифестов \texttt{pr-validation.yml} и \texttt{release.yml}. Пайплайн разделен на два этапа:
	\begin{enumerate}
		\item \textbf{Статический анализ (Linting)}: На данном этапе специализированный инструмент \texttt{ShellCheck} проводит глубокий аудит bash-скриптов на наличие логических ошибок, уязвимостей и соответствие POSIX-стандартам. Дополнительно утилита \texttt{yamllint} верифицирует синтаксис декларативных манифестов развертывания.
		\item \textbf{Динамические тесты}: В изолированном облачном ранере Ubuntu разворачивается точная копия целевого стенда, после чего запускается \texttt{test\_mongo.sh}. Слияние веток кода блокируется до тех пор, пока система не подтвердит успешное прохождение всех тестов.
	\end{enumerate}
	
	\newpage
	
	% ================= ЗАКЛЮЧЕНИЕ =================
	\section*{Заключение}
	\addcontentsline{toc}{section}{Заключение}
	В ходе выполнения данного курсового проекта был успешно спроектирован, оптимизирован и развернут отказоустойчивый распределенный кластер NoSQL СУБД MongoDB в среде операционной системы Linux. Использование концепции Infrastructure as Code и контейнеризации позволило полностью автоматизировать все этапы подготовки окружения, развертывания компонентов и проведения проверочных испытаний.
	
	Реализованный комплекс мер по оптимизации движка WiredTiger исключил риски аварийного завершения процессов базовыми механизмами ядра Linux OOM Killer. Построенная многоуровневая система защиты на базе сетевой изоляции Docker-контура, жестких правил межсетевого экрана UFW и криптографической аутентификации узлов \texttt{keyFile} обеспечила высокий уровень безопасности обрабатываемых данных. Внедрение современных инструментов контроля качества кода и облачного CI/CD пайплайна гарантирует стабильность и воспроизводимость разработанной инфраструктуры на протяжении всего жизненного цикла ее эксплуатации.
	
\end{document}